
\subsubsection{5.1 Concentration Measurement Validity (H-E1, H-M1)}

We successfully computed HHI for all 21 venue-year combinations.

\begin{table}[htbp]
\centering
\resizebox{\columnwidth}{!}{%
\begin{tabular}{ll}
\toprule
Metric & Value \\
\midrule
HHI Range & 0.007 -- 0.046 \\
HHI Variance & 0.00014 \\
Mean HHI & 0.017 \\
Unique task categories & 1,666 \\
\bottomrule
\end{tabular}
}
\end{table}

H-M1 tests whether HHI correlates with simpler concentration metrics. Mann-Whitney U test comparing high-HHI vs. low-HHI groups (split at median HHI = 0.0115) yields:

\begin{table}[htbp]
\centering
\resizebox{\columnwidth}{!}{%
\begin{tabular}{ll}
\toprule
Metric & Value \\
\midrule
Mann-Whitney U & 104.0 \\
p-value & 0.000319 \\
High-HHI mean top-5 share & 0.226 \\
Low-HHI mean top-5 share & 0.147 \\
Spearman $\rho$ (HHI vs. top-5 share) & 0.90 \\
Spearman p-value & < 0.001 \\
\bottomrule
\end{tabular}
}
\end{table}

The high-HHI group shows 54\% higher top-5 share than the low-HHI group. \textbf{Gate: PASSED.}

\subsubsection{5.2 Standard Propagation Mechanism (H-M2, H-M3)}

\textbf{H-M2}: Logistic regression tests whether venue-level concentration in year t-1 predicts individual paper benchmark adoption in year t.

\begin{table*}[htbp]
\centering
\resizebox{\dimexpr 2\columnwidth+\columnsep\relax}{!}{%
\begin{tabular}{lllll}
\toprule
Model & $\beta (prior_HHI)$ & p-value & Odds Ratio & N \\
\midrule
Proposed & 56.75 & 1.5 $\times 10^{-11}$ & 4.4 $\times 10^{24}$ & 10,636 \\
With venue controls & 25.66 & 0.041 & 1.4 $\times 10^{11}$ & 10,636 \\
\bottomrule
\end{tabular}
}
\end{table*}

The extreme magnitude of the odds ratio reflects quasi-complete separation in the data due to the narrow HHI range (0.007--0.046) and should be interpreted as directional evidence of a strong positive relationship. Model fit is poor (Hosmer-Lemeshow $\chi^2$ = 84.5, p < 0.001), suggesting additional predictors may improve calibration. \textbf{Gate: PASSED.}

\textbf{H-M3}: Citation-based benchmark propagation represents the strongest finding.

\begin{table}[htbp]
\centering
\resizebox{\columnwidth}{!}{%
\begin{tabular}{llll}
\toprule
Pair Type & Mean Jaccard & N pairs & Std \\
\midrule
Citing pairs & 0.318 & 77,963 & 0.210 \\
Random pairs & 0.014 & 77,747 & 0.076 \\
\bottomrule
\end{tabular}
}
\end{table}

\begin{table}[htbp]
\centering
\resizebox{\columnwidth}{!}{%
\begin{tabular}{ll}
\toprule
Metric & Value \\
\midrule
Mann-Whitney U & 5.93 $\times 10^9$ \\
p-value & < 0.0001 \\
Cohen's d & 1.93 \\
\bottomrule
\end{tabular}
}
\end{table}

Papers that cite each other (operationalized via task co-occurrence) share benchmarks at approximately 24 times the rate of random pairs. The effect size (d = 1.93) indicates nearly two standard deviations of separation between distributions. \textbf{Gate: PASSED.}

\subsubsection{5.3 Diversity-Concentration Correlation (H-M4)}

H-M4 tests whether low benchmark diversity correlates with high concentration using benchmark breadth (count of distinct task labels per paper) as a proxy for evaluation coverage.

\begin{table}[htbp]
\centering
\resizebox{\columnwidth}{!}{%
\begin{tabular}{ll}
\toprule
Metric & Value \\
\midrule
Spearman $\rho$ & -0.332 \\
p-value & 0.166 \\
N venue-years & 19 \\
\bottomrule
\end{tabular}
}
\end{table}

Per-venue results:

\begin{table}[htbp]
\centering
\resizebox{\columnwidth}{!}{%
\begin{tabular}{llll}
\toprule
Venue & Spearman $\rho$ & p-value & Significant \\
\midrule
NeurIPS & -0.500 & 0.253 & No \\
ICML & -0.900 & 0.037 & Yes \\
ICLR & -0.179 & 0.702 & No \\
\bottomrule
\end{tabular}
}
\end{table}

The overall correlation is in the expected negative direction but fails to reach statistical significance at $\alpha$ = 0.05. Only ICML shows a significant strong negative effect. \textbf{Gate: FAILED.} The relationship between concentration and diversity shows a trend in the expected direction but insufficient evidence for a robust conclusion.

\subsubsection{5.4 Temporal Feedback Test (H-M5)}

The critical test for epistemic lock-in uses panel regression with two-way fixed effects.

\begin{table}[htbp]
\centering
\resizebox{\columnwidth}{!}{%
\begin{tabular}{llll}
\toprule
Parameter & Estimate & 95\% CI & p-value \\
\midrule
$\beta (HHI_{t-1})$ & +1.603 & [-0.37, 3.58] & 0.098 \\
$R^2$ (within) & 0.682 & --- & --- \\
N observations & 18 & --- & --- \\
\bottomrule
\end{tabular}
}
\end{table}

The coefficient is \textbf{positive}, not negative---the opposite direction predicted by lock-in theory. The confidence interval includes zero, and the result is not statistically significant at conventional thresholds.

Robustness checks:

\begin{table}[htbp]
\centering
\resizebox{\columnwidth}{!}{%
\begin{tabular}{llll}
\toprule
Specification & $\beta (HHI)$ & p-value & $R^2$ \\
\midrule
Lag-2 & 0.457 & 0.441 & 0.529 \\
Entity-only FE & 2.935 & < 0.001 & 0.826 \\
Delta spec ($\Delta HHI \rightarrow \Delta Entropy)$ & 0.580 & 0.770 & 0.007 \\
\bottomrule
\end{tabular}
}
\end{table}

Granger causality tests could not be reliably conducted due to insufficient per-venue observations (7 years per venue with maxlag=2 leaves only 4--5 observations per test).

\textbf{Gate: FAILED.} There is no evidence that high concentration in year t-1 predicts lower diversity in year t.

\subsubsection{5.5 Summary of Predictions}

\begin{table}[htbp]
\centering
\resizebox{\columnwidth}{!}{%
\begin{tabular}{llll}
\toprule
Prediction & Expected & Observed & Status \\
\midrule
P1: $HHI_{t-1}$ predicts Entropy\_t decline & $\beta$ < 0 & $\beta$ = +1.60, p = 0.098 & REFUTED \\
P2: HHI Granger-causes Entropy & Significant & 0/3 venues & INCONCLUSIVE \\
P3: High-HHI papers use fewer datasets & Negative correlation & $\rho$ = -0.33, p = 0.17 & PARTIAL \\
\bottomrule
\end{tabular}
}
\end{table}

